\section{Exact finite transfer inputs}\label{sec:transfer}

We use Paper~51's actual positive-Hopf paired-cap surgery, ribbon routing and signed angular transfer \cite[section on actual Wilson insertions]{p51}; see \cite{rt,kir} for the ribbon and punctured modular conventions. A vertex has a zero-framed backbone with caps of framings \(+2,-2\). At level \(r=k+2\), put \(\theta=\pi/r\), \(q=e^{i\theta}\), and let \(V_m\) denote the one-based color of quantum dimension \(\sin(\theta m)/\sin\theta\). The original backbone color is \(r-m_v\); the unpierced negative cap forces each \(m_v\) odd. For fixed integer auxiliary spins, exact simple-current removal supplies one factor \((-1)^r\) per tree edge. With \(N\) vertices the invariant is
\begin{equation}\label{eq:finite}
 Z_{r,W}=(-1)^{r(N-1)}(2/r)^{(N+1)/2}\sin(\theta)^{-(N-1)}
       \sum_{\boldsymbol m\ {\rm odd}}
        \prod_v(1+\cos(\theta m_v))\,\cK_{r,W}(\boldsymbol m).
\end{equation}
All radial or ratio marks multiply this actual summand. The heavy traces in \(\cK_{r,W}\) are normalized; the closed auxiliary trace has its usual categorical normalization. A pierced positive cap on a packet channel of stripped color \(m+2s\) contributes, relative to its unpierced value,
\begin{equation}\label{eq:cap}
 (-1)^s e^{-i\theta(ms+s^2)}
       \frac{\sin(\theta m/2)}{\sin(\theta(m+2s)/2)} .
\end{equation}
The ordinary Hopf transfer is
\begin{equation}\label{eq:lambda}
 \lambda_0(m,n)=\frac{\sin\theta\,\sin(\theta mn)}
                        {\sin(\theta m)\sin(\theta n)} .
\end{equation}

We recall precisely which results of Paper~51 are inputs. Cutting the true tree strips gives a fixed finite sum of products of one-site factors and bounded-rank punctured Hopf transfers, with uniformly bounded suprema and variation for the one-site factors on each parity branch. Fold each original label \(m\) to \(\min(m,r-m)\); reflection signs are fixed on each folded progression. If the largest folded neighbor of a summed label \(n\) is \(M\), the signed conditional sum over any interval is bounded by \(Cr/M\), and over \(Mn/r>K\ge1\) by \(Cr/(MK)\), also with inserted one-site profiles of bounded variation. This is a signed prefix estimate, not an absolute bound on the multiple-color summand.

The other input is the compact-product symbol for the literal routed diagram. Let \(P_W(\boldsymbol u)\) be the product of ordinary characters of the ordered positive-cap half-turn words \(iu_v\cdot\sigma\), \(u_v\in S^2\), with Pauli matrices \(\sigma\). All representatives follow the stipulated exterior corridors with no extra heavy clasp. Then \(\cK_{r,W}\) differs by \(o(1)\) from
\begin{equation}\label{eq:symbol}
 \int P_W(\boldsymbol u)e^{i\pi\sum_{e=vv'}w_e u_v\cdot u_{v'}}
               \prod_v d\Pi(u_v),\qquad w_e=m_v m_{v'}/r,
\end{equation}
uniformly when the \(w_e\) stay bounded, \(\min_v m_v\to\infty\), and \(\max_v m_v/r\to0\). Here \(d\Pi\) is spherical probability; the character product is even in each axis. This theorem concerns the actual normalized diagram, not a diagonal insertion of a function of abstract angles.

\section{Physical ratio gates on matched trees}

\subsection{The resolved angular measure}
Let $G$ be a fixed tree with its unique perfect matching
$M=\{a_jb_j\}_{j=1}^d$, so $N=2d$. An unmatched edge $a_i b_j$
defines an arrow $i\to j$. Choose a subset $E'$ of these arrows,
put $p=d-|E'|$, and contract $E'$ in the matching tree. The remaining
arrows define the order polytope $\Omega_{E'}\subset[0,1]^p$.
Write $\iota_{E'}(u)$ for the full logarithmic array with values
$(u_C,1-u_C)$ at each matched pair in a contracted component $C$.

We use exactly the fixed integer-spin Wilson ribbon representatives
declared in Paper 51, Section ``Actual Wilson insertions and their
joint angular limit'': passages through the designated positive cap
disks, fixed exterior tree corridors, and no new heavy-heavy crossing.
Their joint central symbol is $P_W(a,b)$, a polynomial even under
reversal of each spherical axis. Let $d\Pi$ be spherical probability.
For positive parameters $\rho=(\rho_e)_{e\in E'}$, set
\[
 A_j(a;\rho)=a_j+\sum_{(i\to j)\in E'}\rho_{ij}a_i.
\]
If $A\ne0$, let $d\sigma_A$ be normalized circle measure on
$\{b\in S^2:b\cdot A=0\}$. Define
\begin{equation}\label{eq:ratio-measure}
 d\nu_\rho(a,b)=\prod_j d\Pi(a_j)
       \prod_j\frac{d\sigma_{A_j(a;\rho)}(b_j)}{|A_j(a;\rho)|}.
\end{equation}
The zero sets of the $A_j$ have zero product-sphere measure; the
formula denotes its integrable extension, with no additional atoms.

\begin{lemma}\label{lem:ratio-finite}
The measure $\nu_\rho$ is finite and positive, and
$\nu_\rho(1)\le1$ for every positive $\rho$.
\end{lemma}
\begin{proof}
Choose a sink $j$ of the selected arrow forest. Then $a_j$ appears
only in $A_j$, with coefficient one. For a fixed vector $V$,
direct polar integration gives
\[
 \int_{S^2}\frac{d\Pi(a)}{|a+V|}
   =\frac1{\max(1,|V|)}\le1.
\]
Integrating $a_j$ and deleting its selected arrows bounds the
remaining integral by the same expression on a smaller forest.
Iteration proves the assertion. In particular all possible
$A_j=0$ singularities are integrable.
\end{proof}

\subsection{The quantum transition theorem}
At each vertex choose a fixed smooth profile $g_v(m_v/r)$ equal to
one near zero and supported in $[0,c)$, with $c<1/2$. These ordinary
negative-center profiles are retained. Let $F$ be a fixed smooth
joint test of all physical logarithmic radii $\log m_v/\log r$.
For $e=(i\to j)\in E'$, choose
$\chi_e\in C_c^\infty((0,\infty))$ and insert the literal
finite-level mark $\chi_e(m_{a_i}/m_{a_j})$. No product annulus or
Abel factor is part of this marked invariant.

\begin{theorem}[Physical ratio-resolved Wilson law]
\label{thm:ratio}
For the routes and marks just specified, put
$c_G=\sqrt2\,\pi^{1-N}$. Then
\begin{align}
 Z_{r,W}[F,g,\chi]
 &=(-1)^{r(N-1)}c_G r^{N-3/2}(\log r)^p
 \bigg[
 \int_{\Omega_{E'}}F(\iota_{E'}(u))\,d^pu\notag\\
 &\qquad\times
 \int_{(0,\infty)^{|E'|}}
      \left(\int P_W\,d\nu_\rho\right)
      \prod_{e\in E'}\chi_e(\rho_e)\frac{d\rho_e}{\rho_e}
 +o(1)\bigg].\label{eq:ratio-quantum}
\end{align}
The statement holds along the full integer sequence $r\to\infty$,
with $G,W$ and all profiles fixed. The profiles $\chi_e$ may be complex.
No uniformity as their ratio supports approach zero or infinity is
asserted.
\end{theorem}

\begin{proof}
Apply the finite-rank transfer expansion of Paper 51 and condition on
all $a$ labels. Sum the $b$ labels first, term by term. At $b_j$ let
$M_j$ be the largest neighboring $a$ label. The signed prefix bound
gives $Cr/M_j\le Cr/m_{a_j}$. Joint smooth logarithmic tests are
handled by their anchored mixed-variation representation in rectangular
prefixes, as in Paper 51. The ordinary profiles have uniformly bounded
variation on the retained lower branch.

The ratio marks confine all anchors in a selected component to fixed
positive multiples of one root. Successive leaf sums, or the change
to one root and one edge log ratio, give
\begin{equation}\label{eq:ratio-count}
 \sum_{\boldsymbol m_A\ {\rm odd}}
    \prod_jm_{a_j}^{-1}
    \prod_{(i\to j)\in E'}
       |\chi_{ij}(m_{a_i}/m_{a_j})|
       =O((1+\log r)^p).
\end{equation}
This lower logarithmic power is built into every subsequent bound.

Temporarily put $w_j=m_{a_j}m_{b_j}/r$. On $w_j>K\ge1$ the prefix
threshold is $K M_j/m_{a_j}$, so the conditional tail is bounded by
\[
 \frac{Crm_{a_j}}{KM_j^2}\le\frac{Cr}{Km_{a_j}}.
\]
Keep all other partner sums signed and use \eqref{eq:ratio-count}.
This costs $O(K^{-1})$ after normalization at the power in
\eqref{eq:ratio-quantum}. The estimate is uniform with arbitrary
additional hard interval restrictions on every other matched product:
at fixed anchors these are precisely partner-label intervals in the
prefix lemma. The ratio parameters are integrated against the fixed
$\chi$ in this assertion; a pointwise-in-$\rho$ finite-level estimate
is not needed. Similarly $w_j<\kappa$ contains
$O(\kappa r/m_{a_j})$ labels whenever nonempty. Bounded categorical
factors, with signed sums at the other partners, give $O(\kappa)$.
Thus all matched products may temporarily be kept in $[\kappa,K]$.

An unselected order violation $m_{a_i}/m_{a_j}>r^\epsilon$ gains
$r^{-\epsilon}$ from the $j$-th conditional bound. A width-$\epsilon$
strip about an unselected order wall has harmonic mass
$O(\epsilon(\log r)^p+(\log r)^{p-1})$. Restrict each selected
component root to $[r^\delta,r^{1-\delta}]$; the discarded mass has
the analogous bound with $\delta$. In the retained strict chambers,
all labels tend to infinity and all $m_v/r$ tend to zero. An
unselected product tends to zero, and a selected product is exactly
$\rho_{ij}w_j$. Paper 51's actual joint symbol theorem consequently
gives, uniformly on compact product boxes,
\[
 \mathcal K_{r,W}\longrightarrow
 \int P_W(a,b)\exp\!\left(i\pi\sum_j
               w_j b_j\cdot A_j(a;\rho)\right)
                  \prod_j d\Pi(a_j)d\Pi(b_j).
\]
The retained profiles $g_v$ are then one and $F$ tends to
$F(\iota_{E'}(u))$.

Each odd partner mesh gives $r\,dw_j/(2m_{a_j})$. The anchor harmonic
measure contributes the density
\[
  2^{-d}(\log r)^p d^pu\prod_e d\rho_e/\rho_e;
\]
The log-incidence Jacobian has modulus one because $E'$ is a forest.
The $2d$ cap profiles tend to $2$ each. The exact common scalar
prefactor of the normalized sum in Paper 51 is asymptotic to
\[
 (-1)^{r(N-1)}2^{d+1/2}\pi^{1-2d}r^{d-3/2}.
\]
Thus the coefficient before the $w$ integration is $2^d c_G$.
This also follows from Paper 51's bounded-product overlap theorem,
by setting its selected product $v_{ij}=\rho_{ij}w_j$ and approximating
$\chi_{ij}(v_{ij}/w_j)$ on compact annuli by finite tensor sums.

Since $P_W$ is even in each $b_j$, the exponential can be replaced
by the product of cosines. Passing the interval-uniform bounds to
the compact-product limit gives $O(K^{-1})$ tails on arbitrary
rectangles, after integration against the fixed ratio profiles.
Multivariable summation by parts, or its continuous version, therefore
identifies the ordinary rectangular improper $w$ integral with its
Abel limit. For the Abel integral the following calculation is
pointwise for almost every axis configuration and ratio:
For $A\ne0$,
\[
 \int_0^\infty e^{-\varepsilon w}\cos(\pi ws)\,dw
      =\frac{\varepsilon}{\varepsilon^2+\pi^2s^2}
      \longrightarrow\delta_0(s),\qquad
 \delta(b\cdot A)d\Pi(b)=\frac{d\sigma_A(b)}{2|A|}.
\]
The spherical Abel mass is at most $(2|A|)^{-1}$. Lemma
\ref{lem:ratio-finite} therefore supplies an integrable dominating
function, also after integrating over the fixed compact ratio support.
The full $w$ integral is $2^{-d}\int P_Wd\nu_\rho$. Exhaust the
temporary product boxes and the root and order strips. Multiplication
by $2^dc_G$ proves \eqref{eq:ratio-quantum}.
\end{proof}

\subsection{An explicit nonzero next-logarithm Wilson signal on \texorpdfstring{$P_4$}{P4}}
Take the path $b_1-a_1-b_2-a_2$ and select its unmatched arrow
$a_1\to b_2$. Write $\rho=m_{a_1}/m_{a_2}$ and set
\[
 x=a_1\cdot b_1,\qquad y=a_1\cdot b_2,\qquad z=a_2\cdot b_2.
\]
Under the product sphere law these tree-edge dots are independent
and uniform on $[-1,1]$. Formula \eqref{eq:ratio-measure} is four
times that law restricted by $\delta(x)\delta(z+\rho y)$.
Its projection to $z$, with $a=\min(1,\rho)$, is consequently
\begin{equation}\label{eq:ratio-p4-law}
 x=0,\quad y=-z/\rho,\qquad
 d(\nu_\rho)_z=\frac1{2\rho}\mathbf1_{[-a,a]}(z)\,dz,
 \qquad \nu_\rho(1)=\frac1{\max(1,\rho)}.
\end{equation}
This is a full joint law with a coherent linear constraint.

For the actual spin-one word through caps $a_2,b_2$, the symbol is
$W_z=4z^2-1$. Hence
\[
 H(\rho):=\int(W_z+1)d\nu_\rho
 =\frac{4\min(1,\rho)^3}{3\rho}
 =\begin{cases}4\rho^2/3,&0<\rho\le1,\\
                 4/(3\rho),&\rho\ge1.
   \end{cases}
\]
With $F=1$ and any nonzero, nonnegative
$\chi\in C_c^\infty((0,\infty))$, Theorem~\ref{thm:ratio} gives
\begin{equation}\label{eq:ratio-nonzero}
 Z_{r,W_z}[g,\chi]+Z_r[g,\chi]
 =(-1)^r\frac{\sqrt2}{\pi^3}r^{5/2}\log r
   \left[\int_0^\infty\chi(\rho)H(\rho)\frac{d\rho}{\rho}
           +o(1)\right].
\end{equation}
The coefficient in brackets is strictly positive. Thus a literal
marked quantum observable has a nonzero next-logarithm-order angular
signal outside the principal perpendicular-pair support.

\begin{corollary}[A fixed hard ratio window]\label{cor:ratio-window}
Keep the declared negative-center profiles and restrict the actual
odd labels by $1\le m_{a_1}/m_{a_2}\le2$. With the same Wilson route
and framing in both terms, put
$D_r=Z_{r,W_z}[g,\mathbf1_{[1,2]}]+Z_r[g,\mathbf1_{[1,2]}]$.
Then
\[
 D_r=(-1)^r\frac{2\sqrt2}{3\pi^3}r^{5/2}\log r
                  +o(r^{5/2}\log r),
 \qquad
 Z_r[g,\mathbf1_{[1,2]}]
 =(-1)^r\frac{\sqrt2}{2\pi^3}r^{5/2}\log r
                  +o(r^{5/2}\log r).
\]
\end{corollary}
\begin{proof}
Approximate the window by smooth profiles differing only within
width $\eta$ of its endpoints. Conditional signed partner bounds
reduce their difference, for either Wilson term or the vacuum, to
the anchor harmonic mass of these strips. For $q=1,2$ this obeys
\[
 \sum_{\substack{m,n\ {\rm odd}\\|m/n-q|\le\eta}}
      \frac1{mn}
 \le C\sum_{n<r}\frac1n\left(\eta+\frac1n\right)
 =O(\eta\log r+1).
\]
The estimate also shows that an exact rational-ratio line has
harmonic mass $O(1)$ and therefore no atom at the logarithmically
normalized order. Let $r\to\infty$ and then $\eta\to0$.
Finally
$\int_1^2H(\rho)d\rho/\rho=2/3$ and
$\int_1^2\nu_\rho(1)d\rho/\rho=1/2$.
\end{proof}

The simultaneous spin-one tests $X=4z^2-1$ and $Y=4y^2-1$ further
give, after normalizing $\nu_\rho$, with $b=\min(1,1/\rho)$,
\[
 \mathbb EX=\frac{4a^2}{3}-1,\quad
 \mathbb EY=\frac{4b^2}{3}-1,\quad
 \operatorname{Cov}(X,Y)=\frac{64a^2b^2}{45}>0.
\]
In particular $\mathbb E z^2/\mathbb E y^2=\rho^2$; separate
independent angle marginals would lose this constraint.

\subsection{Every fixed integer spin}
The two-cap spin-$j$ character is $U_{2j}(z)$, with $U_n$ the
Chebyshev polynomial of the second kind. Its weighted coefficient is
\[
 J_j(\rho)=\frac1\rho\int_0^{\min(1,\rho)}U_{2j}(z)\,dz.
\]
For $\rho\ge1$ this is $1/[(2j+1)\rho]$; for $\rho\le1$ it is
an even polynomial in $\rho$. Subtract the principal perpendicular
value by defining
$H_j(\rho)=J_j(\rho)-(-1)^j/\max(1,\rho)$.
Then $H_j=O(\rho^2)$ at zero and $O(\rho^{-1})$ at infinity, and
\begin{equation}\label{eq:ratio-spin-integral}
 \int_0^\infty H_j(\rho)\frac{d\rho}{\rho}
 =\int_0^1\frac{U_{2j}(z)-(-1)^j}{z}\,dz
 =2(-1)^{j+1}\sum_{\substack{1\le k\le j\\k\ {\rm odd}}}\frac1k.
\end{equation}
Indeed absolute convergence permits exchanging the first two
integrals. If $I_j$ is the middle integral, the identity
$U_{2j}+U_{2j-2}=2zU_{2j-1}$ gives
$I_j+I_{j-1}=(1-(-1)^j)/j$, with $I_0=0$.

